\documentclass[11pt]{article}
\usepackage{fullpage, fancyhdr, ifthen, amssymb, amsfonts, amsmath, booktabs}
\usepackage{tabularx}
\usepackage{multicol}
\usepackage[margin=0.25in]{geometry}
\usepackage{upquote}
\usepackage{forest}
\usepackage{ulem}
\usepackage{colortbl}
\pagestyle{empty}

\begin{document}

\begin{table}[h]
\centering
\arrayrulecolor{black!10} % Very light table outline
\begin{tabular}{|p{\textwidth}|}
\hline
\begin{minipage}[t]{\textwidth}
\vspace{0pt}  % This is crucial for top alignment
\noindent \textbf{Instructions:} Please answer the following questions making sure to write your answers legibly on the second page. The code below is a segment of a Python application that provides context for the questions below.
\begin{itemize}
\item Since we are writing Python code, please be careful about spacing/indentation. Leave enough room that it is obvious when you want things indented. 
\item Make sure to write legibly.
\item All questions are worth equal points unless otherwise stated.
\item The quiz is cumulative and may include material from previous weeks.
\end{itemize}
\end{minipage} 
\\
\hline
\end{tabular}
\end{table}

{\color{lightgray}\hrule}
\vspace{1cm}
\underline{\textbf{compute.py}}
\begin{verbatim}
from utils import (
    clean_data,
    format_output,
    validate_params
)

def add_values(x, y):
    return x + y

def multiply_values(x, y):
    return x * y

def triple(n):
    return n * 3

def halve(n):
    return n / 2

def negate(n):
    return -n

if __name__ == '__main__':
    result = multiply_values(4, 7)
    print(f"Result: {result}")
\end{verbatim}


\clearpage

\begin{tabularx}{\textwidth}{l|X}
\textbf{Quiz 4B} &   \textbf{Name: } \\
\hline
\end{tabularx}

\vspace{1cm}

\begin{enumerate}
    \item Please write a function called \texttt{transform\_all} that:
    \begin{itemize}
        \item Takes two arguments: \texttt{x} (a value) and \texttt{operations} (a list of functions)
        \item Applies each function in \texttt{operations} to \texttt{x}
        \item Returns a list of results from applying each function to \texttt{x}
    \end{itemize}
    
    For example: \texttt{transform\_all(6, [triple, halve, negate])} should return \texttt{[18, 3.0, -6]}.
    
\vspace{5.5cm}
    
    \item Please rewrite the function \texttt{multiply\_values} so that when you enter \texttt{print(multiply\_values.\_\_doc\_\_)} in Python, it returns the text: \texttt{Multiplies two values together}. 
    
\vspace{4cm}

    \item Please write a function called \texttt{display\_info} that accepts any number and type of keyword arguments (but no positional arguments) and prints each key-value pair in the format: ``\{key\}: \{value\}" (one per line). If the value is a string, instead print: ``\{key\}: Text".
    
\vspace{5.5cm}
\item {[4 Points]} In class we discussed how code complexity can be broken down into two categories: {\it essential} complexity and {\it accidental} complexity. Please list two sources of accidental code complexity.
    
\end{enumerate}
\end{document}


